\clearpage\thispagestyle{empty}
\pdfbookmark[0]{Quellen zum Bericht über Vesta vom Juni 1807}{vesta1807june-sources}
\begin{center}\textbf{Quellen zum Bericht über Vesta vom Juni 1807}\end{center}
\noindent Diese KI-Übertragung in heutiges Deutsch folgt Franz Xaver von Zachs Bericht vom Juni 1807 in \href{https://gdz.sub.uni-goettingen.de/id/PPN236007467}{Gauß’ \textit{Werken}, Band~VI} (Göttingen, 1874), S.~285--288. Von Zach zitiert Heinrich Wilhelm Matthias Olbers und Carl Friedrich Gauß. Der Erstdruck erschien als Beitrag XLIII, \textit{Fortgesetzte Nachrichten über den neuen von Dr. Olbers in Bremen entdeckten Planeten Vesta}, in der \textit{Monatlichen Correspondenz zur Beförderung der Erd- und Himmelskunde}, Band~XV, \href{https://zs.thulb.uni-jena.de/rsc/viewer/jportal_derivate_00237836/Monatliche_Correspondenz_130168688_15_1807_0597.tif}{S.~590}--600. Der Auszug der Werkausgabe beginnt im letzten Absatz auf \href{https://zs.thulb.uni-jena.de/rsc/viewer/jportal_derivate_00237836/Monatliche_Correspondenz_130168688_15_1807_0600.tif}{S.~593} und endet oben auf \href{https://zs.thulb.uni-jena.de/rsc/viewer/jportal_derivate_00237836/Monatliche_Correspondenz_130168688_15_1807_0606.tif}{S.~599}. Die Zeitschriftendigitalisate stammen von der Thüringer Universitäts- und Landesbibliothek Jena.
\par\medskip
Die Werkausgabe führt die historische Fußnote als Anmerkung~27. Ihre Formulierung \textit{des Herr Verlegers} lautet im Erstdruck auf \href{https://zs.thulb.uni-jena.de/rsc/viewer/jportal_derivate_00237836/Monatliche_Correspondenz_130168688_15_1807_0601.tif}{S.~594} \textit{des Herrn Verlegers}. Die ersten Elemente und ihr Vergleich mit den Beobachtungen stehen auf \href{https://zs.thulb.uni-jena.de/rsc/viewer/jportal_derivate_00237836/Monatliche_Correspondenz_130168688_15_1807_0603.tif}{S.~596}, die zweiten Elemente auf \href{https://zs.thulb.uni-jena.de/rsc/viewer/jportal_derivate_00237836/Monatliche_Correspondenz_130168688_15_1807_0605.tif}{S.~598}. Die Fußnote des Erstdrucks auf S.~598 verweist auf den Eintrag zu 91 Virginis in Beitrag XXIV derselben Zeitschrift, Band~IX (März 1804), \href{https://www.digitale-sammlungen.de/en/view/bsb10538601?page=253}{S.~239}, digitalisiert von der Bayerischen Staatsbibliothek. Dieser Verweis gehört zum Erstdruck. Die sechs editierbaren Altarzeichen folgen den Konturen der Werkausgabe. Übertragung und Satz wurden mit KI erstellt.
\clearpage
